\section{Introducción}

El régimen peruano de metas de gestión de los servicios de agua potable y saneamiento reúne dentro de una misma arquitectura regulatoria indicadores que miden objetos distintos. Algunas metas se refieren a condiciones alcanzadas en la prestación, como la cobertura, la continuidad o la presión; otras verifican la ejecución de inversiones, actividades, utilización de recursos o capacidades de gestión. Aunque todos estos objetos son relevantes para la regulación, difieren en su naturaleza económica, horizonte temporal y grado de controlabilidad: la ejecución de una intervención, su entrada en operación y la materialización de sus efectos sobre el servicio no necesariamente ocurren en el mismo periodo. Esta heterogeneidad plantea una cuestión de diseño: si todos estos objetos deben incidir de la misma manera en una medida agregada de cumplimiento o si corresponde asociarlos con instrumentos diferenciados según la función que cumplen.

La agregación de resultados y variables de implementación dentro de un mismo índice puede dificultar su interpretación. Cuando parte de la ponderación se asigna a inversiones, actividades o gasto, disminuye la sensibilidad de la medida frente a las variaciones en los resultados del servicio y el cumplimiento de medios puede compensar un desempeño insuficiente en las condiciones efectivamente alcanzadas. Esta preocupación se relaciona con la literatura sobre incentivos multitarea y congruencia de las medidas de desempeño: cuando el agente realiza varias tareas o la medida utilizada no se encuentra adecuadamente alineada con el objetivo del principal, los incentivos pueden distorsionar la asignación del esfuerzo y favorecer mejoras de la métrica que no se traducen proporcionalmente en el objetivo subyacente \citep{holmstrom1991multitask,baker1992incentive}. En el sector de agua potable, \citet{PicazoTadeo2008Quality} muestran, además, que omitir una dimensión sustantiva como la calidad puede modificar significativamente la evaluación de la eficiencia de las empresas. Esta cuestión de medición tiene, a su vez, una dimensión institucional: en el Perú, el cumplimiento de las metas se evalúa mediante el Índice de Cumplimiento Global (ICG) y los índices de cumplimiento individual (ICI), cuyos resultados se encuentran vinculados con la fiscalización y el régimen sancionador.

La pregunta adquiere especial relevancia tras el cambio introducido por el Decreto Legislativo N.\textdegree~1620. En el diseño anterior, el cumplimiento de las metas de gestión podía condicionar la aplicación total o parcial de los incrementos tarifarios base; la nueva regla desvinculó dichos incrementos del cumplimiento de condiciones, sin eliminar la exigibilidad de las metas. Al modificarse este vínculo tarifario, cobra mayor importancia definir qué información debe transmitir el ICG y qué obligaciones requieren mecanismos específicos de seguimiento, fiscalización y sanción. Este artículo estudia, en ese contexto, si el sistema peruano de metas de gestión puede mejorarse mediante una separación funcional de instrumentos y examina las condiciones para transitar hacia una arquitectura de metas de desempeño. Bajo la alternativa evaluada, el ICG se concentraría en resultados observables de cobertura y calidad del servicio, mientras que las inversiones, actividades, capacidades de gestión y demás dimensiones regulatorias permanecerían sujetas a los instrumentos que correspondan a su naturaleza. La separación no reduce la importancia ni la exigibilidad de estas obligaciones, sino que distingue la medición de las condiciones alcanzadas de la verificación de los medios utilizados para producirlas.

El análisis combina la sistematización del marco regulatorio, evidencia descriptiva de las fiscalizaciones realizadas por la Sunass, experiencia regulatoria comparada y una formalización de los incentivos generados por distintos arreglos institucionales. Para identificar separadamente sus efectos, el modelo compara una secuencia de esquemas que delimita progresivamente el ámbito de la medición financiera, elimina el umbral agregado como condición de activación de las consecuencias, sustituye la verificación financiera de determinadas intervenciones por indicadores de cumplimiento físico o funcional y, finalmente, concentra el ICG en resultados. La secuencia \(A_0\to A\to A'\to A''\to B\) permite aislar cada modificación, en lugar de atribuir a una única reforma diferencias que podrían provenir de varios cambios simultáneos. Una extensión \(B^{+}\) incorpora, de manera separada, un mecanismo acotado de reconocimiento del sobrecumplimiento.

La principal contribución consiste en distinguir tres funciones regulatorias que pueden aparecer superpuestas dentro de un mismo esquema: la verificación de las inversiones, la aplicación de consecuencias sancionadoras y la construcción de la señal agregada de desempeño. La comparación secuencial permite examinar por separado los canales de verificación, de activación sancionadora y de composición del índice, e identificar si y bajo qué condiciones la superposición entre estas funciones modifica los incentivos de ejecución, eficiencia y desempeño. Sobre esta misma arquitectura, la extensión por sobrecumplimiento permite analizar las condiciones bajo las cuales podría reconocerse una mejora adicional sin alterar la medición ordinaria del cumplimiento ni introducir compensaciones entre metas.

La evidencia empírica cumple una función descriptiva y contrafactual acotada. Los resultados de las fiscalizaciones permiten caracterizar diferencias en los niveles de cumplimiento, pero no identifican causalmente las razones que las explican. Asimismo, la simulación contrafactual mantiene constantes los índices de cumplimiento individual y modifica únicamente los componentes incorporados al ICG, por lo que identifica el efecto aritmético de su composición y no la respuesta conductual de las empresas ni el efecto causal de una reforma sobre las inversiones o los resultados del servicio. El modelo constituye, por su parte, una representación estilizada de los incentivos y no determina un diseño regulatorio óptimo. En consecuencia, el análisis no postula la superioridad incondicional de una arquitectura: la conveniencia de separar instrumentos depende de que la fiscalización y el régimen sancionador preserven incentivos suficientes para ejecutar las obligaciones de implementación, de que los indicadores incorporados al ICG representen adecuadamente las condiciones que se pretende hacer visibles, de la controlabilidad de dichas variables y de que la señal de desempeño alcance suficiente difusión y legibilidad. En el caso del sobrecumplimiento, debe considerarse adicionalmente el riesgo de desplazar esfuerzos y recursos desde obligaciones pendientes hacia dimensiones susceptibles de reconocimiento, posibilidad que queda fuera de los supuestos de capacidad gerencial no vinculante utilizados en el modelo.

El documento se organiza como sigue. La Sección~2 presenta el marco regulatorio aplicable y la evidencia descriptiva sobre el cumplimiento de las metas de gestión. La Sección~3 examina si la heterogeneidad de los objetos comprendidos en dichas metas justifica un tratamiento diferenciado y desarrolla el principio de separación de instrumentos a partir de su función regulatoria. La Sección~4 formaliza los incentivos asociados con distintos arreglos institucionales y compara secuencialmente los esquemas \(A_0\), \(A\), \(A'\), \(A''\), \(B\) y la extensión \(B^{+}\). La Sección~5 discute conjuntamente los resultados del análisis, examina sus implicancias para la composición y función del ICG, la fiscalización y el régimen sancionador, y desarrolla las condiciones que requeriría una eventual implementación de los cambios evaluados, incluido el tratamiento del sobrecumplimiento. Finalmente, la Sección~6 presenta las conclusiones. El Anexo~\ref{app:calculo_sancionador} desarrolla, como complemento técnico, la correspondencia entre el incumplimiento de los índices, las inversiones no ejecutadas y la base económica utilizada para determinar la sanción, mientras que el Anexo~\ref{anexo:experiencia_internacional} presenta la experiencia regulatoria internacional utilizada como referencia comparativa.
